\chapter{Conclusion}
\label{chapter:conclusion}

\section{Future Work}
One promising direction for future work is to extend LeanCat to support more complex model comparisons, such as between RISC-V and TSO, or between LKMM and the underlying architectures it is designed to support. As we discussed in Section~\ref{sec:limitations}, the Cat language leaves many properties implicit because herdtools7 handles them computationally. For example, Cat assumes coherence order is a strict total order per location without stating this formally. In Lean 4, every such assumption must be made explicit as a wellformedness condition before proofs become possible. For simple models like SC and TSO, the number of such implicit assumptions is manageable. However, for more complex architectures like RISC-V, the number of implicit assumptions grows significantly, and identifying and formalizing them correctly requires substantial additional effort. LeanCat already provides the foundational infrastructure for such comparisons --- the \texttt{defcat} command, the \texttt{acyclicMono} lemma, and the candidate execution framework are all in place. The remaining work is to identify the implicit assumptions for each new model and add them as wellformedness conditions, which is a well-defined engineering task rather than a fundamental research challenge. We therefore expect that extending LeanCat to support RISC-V and other architectures is achievable in future work, building directly on the infrastructure developed in this thesis.